\reportelement{Введение}
Цель работы --- подключить PostgreSQL к персональному ассистенту на FastAPI, организовать хранение заметок в именованном томе и проверить взаимодействие контейнеров в пользовательской сети Docker. Я продолжил работу с ветки ЛР № 2 и создал ветку \texttt{lab-3}. Проверки выполнены в Docker Desktop на macOS.

\begin{table}[H]
\caption{Персональные параметры лабораторной работы}
\begin{tabularx}{\textwidth}{|p{61mm}|X|}
\hline
Параметр & Значение \\\hline
ФИО & Ивановский Роман Евгеньевич \\\hline
Группа & К3320 \\\hline
Номер ИСУ & 466021 \\\hline
Порт PostgreSQL на хосте & 5430 + 21 = 5451 \\\hline
Порт ассистента на хосте & 8000 + 21 = 8021 \\\hline
Префикс & ivanovsky \\\hline
Учебный ключ ассистента & secret21ivanovsky \\\hline
Имя базы данных & ivanovsky21 \\\hline
\end{tabularx}
\end{table}

На первом этапе я использовал порт \texttt{5451} для доступа к PostgreSQL через хост. После переноса в пользовательскую сеть PostgreSQL принимал соединения на внутреннем порту \texttt{5432}. Ассистент на всех этапах был доступен на порту \texttt{8021}. Учебный ключ я передавал в контейнер через переменную \texttt{APP\_SECRET}.

\chapter{Именованный том и подключение PostgreSQL}
Я создал том \texttt{ivanovsky-pgdata} и подключил его к каталогу \texttt{/var/lib/postgresql/data} контейнера PostgreSQL 15. Для связи с ассистентом на первом этапе я опубликовал порт базы \texttt{5451} и указал адрес \texttt{host.docker.internal:5451} в переменной \texttt{DATABASE\_URL}.

\begin{code}
docker volume create ivanovsky-pgdata
docker build -f Dockerfile.optimized -t ivanovsky-assistant:v3 .
\end{code}

При запуске FastAPI создаёт таблицу \texttt{notes}, если она ещё отсутствует. Эндпоинты \texttt{POST /notes} и \texttt{GET /notes} записывают и читают данные через PostgreSQL. В обработчике \texttt{/health} я добавил запрос \texttt{SELECT 1}.

На macOS порт \texttt{8021} был занят службой \texttt{com.apple.ftp-proxy}. После временной выгрузки службы я запустил ассистента на персональном порту. Запросы к \texttt{/} и \texttt{/health} завершились успешно. Я создал заметку «ЛР3: данные в БД» и получил её запросом \texttt{GET /notes}.

\reportscreenshot{images/step-1-5.png}{Проверка HTTP-запросами и сохранение первой заметки}{fig:stage-1}

\chapter{Проверка сохранности заметки}
Я остановил и удалил контейнер ассистента, затем создал новый контейнер из того же образа. PostgreSQL и именованный том продолжали работать. После запуска нового контейнера я запросил список заметок.

\begin{code}
docker stop ivanovsky-assistant
docker rm ivanovsky-assistant
curl -4 -fsS http://localhost:8021/notes |
  python3 -m json.tool --no-ensure-ascii
\end{code}

В ответе сохранилась заметка «ЛР3: данные в БД», созданная до пересоздания ассистента. На контрольном снимке также видны время проверки, логин и имя хоста.

\reportscreenshot{images/step-2-4.png}{Получение исходной заметки после пересоздания ассистента}{fig:stage-2}

\chapter{Подключение файла конфигурации}
Я создал файл \texttt{config/app.conf} с параметрами ассистента и значением \texttt{owner=ivanovsky}. При запуске контейнера каталог \texttt{config} с хоста я подключил к \texttt{/app/config} в режиме только для чтения с параметром \texttt{-v "\$PWD/config:/app/config:ro"}.

\begin{code}
docker exec ivanovsky-assistant-config cat /app/config/app.conf
docker restart ivanovsky-assistant-config
docker exec ivanovsky-assistant-config cat /app/config/app.conf
\end{code}

Между двумя проверками я изменил на хосте значение \texttt{maintenance\_mode} с \texttt{false} на \texttt{true}. После перезапуска новое значение отобразилось внутри контейнера без пересборки образа. Проверка подтверждает доступ контейнера к обновлённому файлу.

\reportscreenshot{images/step-3-3.png}{Файл конфигурации внутри контейнера до и после изменения на хосте}{fig:stage-3}

\chapter{Пользовательская сеть и изоляция базы}
Я создал сеть \texttt{ivanovsky-backend} и пересоздал в ней PostgreSQL и ассистента. Для PostgreSQL я сохранил том \texttt{ivanovsky-pgdata} и убрал публикацию порта на хосте. Ассистент подключался к базе по адресу \texttt{ivanovsky-postgres:5432} внутри сети.

После переноса \texttt{/health} подтвердил подключение к PostgreSQL. Я создал заметку «ЛР3: сеть работает!» и получил обе заметки запросом \texttt{GET /notes}.

В стандартной сети контейнер Alpine получил ошибку разрешения имени \texttt{ivanovsky-postgres}. В сети \texttt{ivanovsky-backend} имя разрешилось, и \texttt{ping} получил два ответа без потерь.

\reportscreenshot{images/step-4-4-outside.png}{Ошибка разрешения имени вне сети ivanovsky-backend}{fig:stage-4}

\reportscreenshot{images/step-4-4-inside.png}{Проверка доступности контейнера внутри сети ivanovsky-backend}{fig:stage-5}

\chapter{Итоговая проверка инфраструктуры}
В финальном выводе от 24.11.2026 я проверил состояние обоих контейнеров, именованного тома и пользовательской сети. Ассистент имел статус \texttt{healthy} и проброс порта \texttt{8021}. В строке PostgreSQL отображался только внутренний порт \texttt{5432/tcp}.

\begin{code}
docker ps --filter 'name=ivanovsky-'
docker volume ls --filter 'name=ivanovsky-pgdata'
docker network ls --filter 'name=ivanovsky-backend'
\end{code}

\reportscreenshot{images/step-5-1.png}{Итоговое состояние контейнеров, тома и сети}{fig:stage-6}

Репозиторий работы, ветка lab-3.\ \url{https://github.com/st1verka/course-cont-orch-app/tree/lab-3}

\reportelement{Заключение}
Я подключил PostgreSQL к ассистенту на FastAPI и проверил, что заметка сохраняется после пересоздания контейнера приложения. Подключённый через bind mount конфигурационный файл обновился внутри контейнера после изменения на хосте. В итоговой схеме ассистент подключается к базе по имени контейнера в пользовательской сети, а на хосте опубликован только порт ассистента.

\begin{thebibliography}{9}
\bibitem{source1} Курс «Контейнеризация и оркестрация приложений». Задание лабораторной работы № 3.\ \url{https://github.com/itmo-faculty-applied-informatics/course-cont-orch-app/blob/main/labs/03-persistence.md}
\bibitem{source2} Репозиторий с результатами лабораторной работы, ветка lab-3.\ \url{https://github.com/st1verka/course-cont-orch-app/tree/lab-3}
\end{thebibliography}
